\section{Заключение}

После выполнения лабораторной работы у меня появились ответы на следующие вопросы.

\paragraph{Почему используемый критерий останова может быть неточным?} 
В работе мы допустили, что метод сходится быстро. Если же сходимость «медленная» ($q \to 1$), то разность двух соседних приближений может быть очень малой, в то время как само решение всё еще находится далеко от текущей точки. В таких ситуациях использование упрощенного критерия без учета $q$ может привести к преждевременной остановке алгоритма до достижения требуемой точности.

\paragraph{Почему ряд тестов не получилось пройти на платформе Code\&Test?} 
Полагаю, что причина кроется в высокой чувствительности к начальным условиям и виду эквивалентной системы. Если для тестовой системы спектральный радиус матрицы Якоби $q \ge 1$ в окрестности корня, то итерационный процесс расходится.

\paragraph{Какая асимптотическая сложность у алгоритма?}
Временная сложность метода простых итераций составляет $\mathcal{O}(M \cdot n)$, где $M$ — количество затраченных итераций, а $n$ — размерность системы. Как следует из блок-схемы, на каждой итерации происходит однократное вычисление вектора-функции $\Phi(x)$ и расчет нормы невязки за линейное время $\mathcal{O}(n)$. Здесь мы допускаем, что расчёт значения отображения \( \Phi \) занимает константное время.

Пространственную сложность можно оценить $\mathcal{O}(n)$, поскольку алгоритм требует выделения памяти лишь под два вектора размерности $n$, что эффективнее прямых методов, требующих хранения квадратных матриц.

\paragraph{В чём преимущество метода простых итераций?} 
Если сравнивать с методом Ньютона, главное преимущество --- вычислительная простота одного шага. Метод простых итераций не требует вычисления элементов матрицы Якоби и решения системы линейных алгебраических уравнений на каждой итерации. Если сравнивать с методом градиентного спуска, то простому итерационному процессу не нужно подбирать какие-либо гиперпараметры \textit{(в частности, шаг обучения)}.

\paragraph{Какие недостатки обнаружены у метода простых итераций?} 
Главный недостаток --- линейная скорость сходимости, которая существенно уступает квадратичной сходимости метода Ньютона.

Кроме того, метод требует трудоемкого предварительного этапа: ручного преобразования исходной системы $F(x) = 0$ к эквивалентному виду $x = \Phi(x)$. При этом преобразование должно быть выполнено так, чтобы гарантировалось выполнение достаточного условия сходимости.

Конечно, есть эмпирические модификации метода простых итераций, которые позволяют формировать преобразование \( \Phi \) автоматически. Однако их рассмотрение выходит за рамки лабораторной работы и лекционного материала.

\paragraph{Почему число итераций у трансцендетной системы резко возросло?} 

В разделе с условиями сходимости метода было сказано, что её скорость линейна и зависит от нормы матрицы Якоби. Увеличение количества итераций лишь говорит о том, что \( q \) оказалось очень близко к единице. То есть каждый последующий шаг лишь незначительно приближает нас к ответу, и алгоритму требуется большое число итераций для достижения требуемой точности.